% Active elementary proof from numbered source 2615.
\subsection{Elementary local capacity bounds}
Write $f(a,b)=1-c_{\max}(a,b)$ for $a,b\ge0$, $a+b\le1$, and put
$m=\min(a,b)$. Let $h=\sqrt3/2$ and let $c(m)\in[h,1]$ be the root of
\[
 F(C,m)=C^4-C^2+mC-m^2=0,\qquad \ell(m)=1-c(m).
\]
It is unique because $F_C>0$ on $[h,1]$,
$F(h,m)=-(m-h/2)^2\le0$, and $F(1,m)=m(1-m)\ge0$.
Define
\[
 q(m)=\frac{m(1-m)}2,\qquad
 \beta(m)=\frac{2m}{4+3m},\qquad \kappa(m)=1+\frac32m.
\]
\begin{lemma}[Own-ray deficit bounds]
\label{lem:bc-capacity-tools}
For $a,b\ge0$, $a+b\le1$, and $m=\min(a,b)$,
\[
 q(m)\le\ell(m)\le f(a,b)\le m.
\]
For $m\le1/4$, also $\beta(m)\le\ell(m)$ and $\beta(m)\ge q(m)$.
\end{lemma}
\begin{proof}
The upper bound and capacity antitonicity follow from
\zcref{lem:new-common-pair-domination}. By the selected local formula
\zcref{prop:new-exact-local-set}, for $m\le h/2$ the diagonal pair lies
in the quartic cell (its cell polynomial is $m^2(16m^2-3)\le0$), so
$c_{\max}(a,b)\le c_{\max}(m,m)=c(m)$. For $m\ge h/2$,
antitonicity gives $c_{\max}(a,b)\le c_{\max}(h/2,h/2)=h\le c(m)$.
Thus $f(a,b)\ge\ell(m)$.

For $c=c(m)$, the root equation and a factorization give
\[
 \ell c^2(1+c)=m(c-m),
\]
\[
 2(c-m)-(1-m)c^2(1+c)
 =(1-c)\bigl((1-m)c(c+2)-2m\bigr)\ge0.
\]
The bracket is at least
$\tfrac12\cdot\tfrac34\cdot\tfrac{11}4-1=1/32$, since $c\ge h>3/4$
and $m\le1/2$. Division by $c^2(1+c)>0$ proves $\ell\ge q(m)$.

For the stronger bound on $m\le1/4$, put
\[
 B(m)=81m^4+405m^3+656m^2+272m-128.
\]
Direct substitution gives
\[
 F(1-\beta(m),m)=-\frac{m^2B(m)}{(4+3m)^4}.
\]
The polynomial $B$ is increasing and $B(1/4)=-3163/256<0$.
Also $1-\beta(m)\ge17/19>h$. Monotonicity of $F$ in $C$ implies
$\beta(m)\le\ell(m)$. Finally
$\beta(m)-q(m)=m^2(3m+1)/(2(3m+4))\ge0$. At $m=0$ all bounds are immediate.
\end{proof}

\begin{lemma}[Quarter-range slack envelope]
\label{lem:bc-slack-envelope}
For $a,b\ge0$, $a+b\le1$, $m=\min(a,b)\le1/4$, and $\delta=1-a-b$,
\[
 f(a,b)\ge\max\{\beta(m),\ m-\kappa(m)\delta\}.
\]
The two entries meet at $\delta=\beta(m)$. This estimate is asserted only
on $m\le1/4$; the broader historical envelope is not needed here.
\end{lemma}
\begin{proof}
Since $m-\kappa(m)\beta(m)=\beta(m)$, the baseline in
\zcref{lem:bc-capacity-tools} proves the claim for $\delta\ge\beta(m)$.
The case $m=0$ is immediate. Suppose $m>0$ and $0<\delta\le\beta(m)$.
Put $M=1-m-\delta$. Then $M\ge m$ (indeed
$1-2m-\beta(m)\ge15/38>0$) and $1-\delta\ge1-\beta(m)\ge c(m)$.
The selected triangular branch of \zcref{prop:new-exact-local-set}
therefore applies: the capacity is the smaller root of
\[
 H(C)=\bigl((1-\delta)^2-1\bigr)C^2+MC-M^2.
\]
At $\overline C=1-m+\kappa(m)\delta$, expansion gives
$H(\overline C)=\delta N_m(\delta)/4$, where
\[
\begin{split}
 N_m(\delta)={}&(3m+2)^2\delta^3-10m(3m+2)\delta^2\\
 &+(28m^2-22m-20)\delta+14m(1-m).
\end{split}
\]
Because $\delta\le\beta(m)\le m/2\le1/8$,
\[
 N_m'(\delta)\le3(11/4)^2(1/8)^2+7/4-20=-18325/1024<0.
\]
At the right endpoint,
\[
 N_m(\beta(m))=-\frac{2mB(m)}{(4+3m)^3}>0.
\]
Thus $H(\overline C)>0$. The quadratic is concave and the capacity is
its selected smaller root, so $c_{\max}(M,m)<\overline C$. This proves
the affine lower bound. At $\delta=0$, the exact capacity is $1-m$.
\end{proof}

\begin{lemma}[Coupled nonuniform capacity inequality]
\label{lem:bc-coupled-capacity}
If $0<x\le y<1$ and $x/2\le z\le y$, put
\[
 r=f(1-y,z),\qquad t=f(1-z,x/2).
\]
Then
\[
 2(r+t)>(1-y+r)(2x-y+r).
\]
\end{lemma}
\begin{proof}
Put $u=x/2$ and
\[
 P(y,r,t)=2(r+t)-(1-y+r)(4u-y+r).
\]
The domain is $0<u$, $2u\le y<1$, $u\le z\le y$ and
$0<r,t\le1/2$. By \zcref{lem:bc-capacity-tools},
\[
 P_r=1-4u+2y-2r\ge1-2r\ge0,\qquad P_t=2.
\]
Hence valid lower bounds on either radius may be substituted. These are
scalar comparisons; the exact witness radii are not changed.

\readerheading{Easy regions.}
First let $y\ge1/2$ and put $a=1-y$. Then $0<a\le1/2$,
$u\le(1-a)/2$, and $u\le z\le1-a$.
If $z\ge a$ and $z\le1-u$, substitute $q(a),q(u)$.
The expression is concave in $u$; at $u=0$ it equals
$2q(a)+(a+q(a))(1-a-q(a))>0$, and at $u=(1-a)/2$ it equals
$(1-a)^3(1+a)/4>0$.
If $z\ge a$ and $z\ge1-u$, both radii are at least $q(a)$ because
$a\le1-z\le u\le1/2$ and $q$ is increasing. The substituted expression
decreases with $u$ and at $u=(1-a)/2$ equals
$a(1-a)(a^2-a+2)/4>0$.
If $z\le a$, both radii are at least $q(u)$. The expression is concave
in $a\in[u,\min(1/2,1-2u)]$, with endpoint values
\[
\begin{array}{c|c|c}
a&P\text{ after substitution}&u\text{ range}\\\hline
u&u(14-43u+14u^2-u^3)/4&(0,1/3]\\
1/2&(1-17u^2+10u^3-u^4)/4&(0,1/4]\\
1-2u&u(-u^3+2u^2+9u-2)/4&[1/4,1/3].
\end{array}
\]
The three brackets respectively decrease to $32/27$, decrease to
$23/256$, and increase from $23/64$, so all are positive.

Now let $y\le1/2$ and $z\ge2u$. Substitute $q(z),q(u)$.
The result increases with $y$, so set $y=z$; it is then concave in
$u\in[0,z/2]$. The value at zero is positive, and the value at
$u=z/2$ is $z^3(2-z)/4>0$.
Finally, for $y\le1/2$ and $1/4\le z\le2u$, the same substitution
allows $y=2u$, giving
\[
 P\ge q(z)(1-q(z))-u+3u^2
 \ge\frac3{32}\frac{29}{32}-\frac1{12}=\frac5{3072}>0.
\]
Here $q(z)(1-q(z))$ increases on $[1/4,1/2]$ and
$-u+3u^2=3(u-1/6)^2-1/12$.

\readerheading{The remaining region and its switches.}
It remains to treat
\[
 0<u\le1/4,\quad u\le z\le\min(2u,1/4),\quad 2u\le y\le1/2.
\]
The smaller capacity coordinates are $z$ and $u$. By
\zcref{lem:bc-slack-envelope} compare with
\[
 R_0=\max\{\beta(z),z-\kappa(z)(y-z)\},\qquad
 T_0=\max\{\beta(u),u-\kappa(u)(z-u)\}.
\]
Put $g(v)=v+\beta(v)$, an increasing rational function. The switches are
$y=g(z)$ and $z=g(u)$.
For fixed $u,z$, the affine branch of $R_0$ gives
\[
 \frac{d^2}{dy^2}P(y,R_0,T_0)=-2(1+\kappa(z))^2<0.
\]
The baseline branch gives $dP/dy=1+4u+2\beta(z)-2y>0$.
Since $g(z)\le g(1/4)=27/76<1/2$, only $y=2u$ or $y=g(z)$
needs to be checked, whenever the latter belongs to the domain.

On $y=2u$,
\[
 P=R_0(1-R_0)+2T_0-2u(1-2u).
\]
Split the $z$ interval at $g(z)=2u$ and $z=g(u)$.
The function $\beta(z)(1-\beta(z))$ is concave, since $\beta$ is concave
and $v(1-v)$ is increasing and concave on $[0,1/2]$.
On the other branch,
\[
 R=2z+\tfrac32z^2-2u-3uz,\qquad R'=2+3z-3u\ge2,\qquad R''=3.
\]
Where it is active, $0<R\le z\le1/4$, so
\[
 (R(1-R))''=3(1-2R)-2(R')^2\le-5<0.
\]
As $T_0$ is affine or constant on each piece, $P$ is piecewise concave.
At $z=u$ it equals $\beta(u)(1-\beta(u))+4u^2>0$; at $z=2u$ it
is $2T_0>0$. The endpoint $z=1/4$ is covered by the easy-region bound,
since the comparison radii dominate the $q$ bounds. Only the switches remain.

On $y=g(z)$,
\[
 P=2\beta(z)+2T_0-(1-z)(4u-z),\qquad z/2\le u\le g(z)/2.
\]
Its $u$ derivative on the baseline branch is
$16/(4+3u)^2-4(1-z)\le-2<0$; on the other branch it is $6u+z>0$.
Thus the minimum lies at an endpoint or $z=g(u)$. The endpoints are
$z=2u$, already checked, and transition A below. The ordering of switches
is not assumed; coincident switches are included.

\readerheading{Three transition checks.}
In A, $y=2u=g(z)$. Using the affine entry as a lower bound for $T_0$ gives
\[
 P\ge\frac{z^2(9z^2+36z+44)}{4(3z+4)^2}>0.
\]
In B, $y=2u$ and $z=g(u)$, so $T_0=\beta(u)$. If $R_B$ denotes the
affine entry of $R_0$, then
\[
 R_B-u(1-2u)=\frac{u^2(9u^2+6u+4)}{2(3u+4)^2}>0.
\]
Since $v(1-v)$ increases on $[0,1/2]$,
\[
 P\ge\frac{u^2(1+19u-4u^2-12u^3)}{3u+4}>0,
\]
where $4u^2+12u^3\le7/16$.
In C, $y=g(z)$ and $z=g(u)$, direct substitution gives
\[
 P=\frac{3u^2(81u^4+441u^3+768u^2+460u+48)}
 {(3u+2)(3u+4)^2(3u+8)}>0.
\]
These regions and transitions are exhaustive. Monotonicity in the radii
therefore proves the inequality for the exact capacities.
\end{proof}
